\section{Mathematical Preliminaries}
\label{sec:preliminaries}

This section establishes the mathematical framework for the Squiggly Line Conjecture and defines key concepts used throughout the paper.

\subsection{Problem Setup}

Let $\mathcal{D} = \{(x_i, y_i)\}_{i=1}^n$ be a set of points in $[-1, 1]^d \times \mathbb{R}$.

For a function $f: [-1, 1]^d \to \mathbb{R}$, define the \textbf{average distance to the dataset}:
$$L(f, \mathcal{D}) = \frac{1}{n} \sum_{i=1}^n \min_{x \in [-1,1]^d} \|(x_i, y_i) - (x, f(x))\|_2$$

\subsection{The Squiggly Line Conjecture}

\begin{definition}[The Squiggly Line Conjecture]
For randomly initialized functions $f_{\theta_1}$ and $f_{\theta_2}$ where $|\theta_1| < |\theta_2|$ (parameter count):
$$\mathbb{E}[L(f_{\theta_2}, \mathcal{D})] < \mathbb{E}[L(f_{\theta_1}, \mathcal{D})]$$
\end{definition}

\subsection{Spatial Coverage Mechanism}

The conjecture relies on three key geometric concepts:

\begin{definition}[Curve Length]
For a function $f: [-1,1]^d \to \mathbb{R}$, its curve length is:
$$\ell(f) = \int_{[-1,1]^d} \sqrt{1 + \|\nabla f(x)\|^2} dx$$
\end{definition}

\begin{definition}[Coverage Volume]
The $\epsilon$-neighborhood of function $f$ is:
$$V_\epsilon(f) = \text{Vol}\{(x, y) : \exists t \in [-1,1]^d, \|(x, y) - (t, f(t))\| < \epsilon\}$$
\end{definition}

\begin{definition}[Expected Minimum Distance]
For uniform random points, the expected minimum distance satisfies:
$$\mathbb{E}[\text{dist}(p, f)] \approx \left(\frac{V_{\text{uncovered}}}{V_{\text{total}}}\right)^{1/d}$$
\end{definition}

\subsection{Polynomial Spaces}

Let $P_d$ be the space of polynomials of degree $d$:
$$P_d = \left\{p(x) = \sum_{k=0}^d a_k x^k : a_k \in \mathbb{R}\right\}$$

For a polynomial $p \in P_d$ with random coefficients $a_k \sim \text{Uniform}(-\alpha, \alpha)$, we can compute:

\textbf{Variance of the polynomial}:
$$\text{Var}[p(x)] = \sum_{k=0}^d \text{Var}[a_k] x^{2k} = \frac{\alpha^2}{3} \sum_{k=0}^d x^{2k}$$

For $x \in [-1, 1]$: $\text{Var}[p(x)] \leq \frac{\alpha^2}{3}(d + 1)$

\textbf{Expected number of local extrema}:
By Kac's theorem, a random polynomial of degree $d$ has expected number of real roots:
$$\mathbb{E}[\#\text{extrema}] \approx \frac{2}{\pi} \log(d) + O(1)$$